\documentclass[11pt,oneside]{article}

\usepackage[a4paper,margin=27mm,headheight=15pt]{geometry}
\usepackage[T1]{fontenc}
\usepackage[utf8]{inputenc}
\IfFileExists{libertinus.sty}{\usepackage{libertinus}}{\usepackage{lmodern}}
\usepackage{microtype}
\usepackage{amsmath,amssymb,amsthm,mathtools,mathrsfs}
\usepackage{bm}
\usepackage{booktabs,longtable,array,tabularx}
\usepackage{graphicx}
\usepackage{pgfplots}
\usepackage{xcolor}
\usepackage{enumitem}
\usepackage{fancyhdr}
\usepackage{titlesec}
\usepackage{listings}
\usepackage{xurl}
\usepackage{hyperref}
\usepackage[nameinlink,noabbrev]{cleveref}

\definecolor{linkblue}{RGB}{25,75,135}
\definecolor{shadegray}{RGB}{246,247,249}
\definecolor{rulegray}{RGB}{195,201,208}
\pgfplotsset{compat=1.18}
\hypersetup{
  colorlinks=true,
  linkcolor=linkblue,
  citecolor=linkblue,
  urlcolor=linkblue,
  pdftitle={Signed Thue--Morse Prefix Sums and the Fabius Function},
  pdfsubject={Sharp convergence rates for discrete Fabius approximants},
  pdfkeywords={Fabius function, Rvachev up-function, Thue-Morse, prefix sums, convergence rate, Fourier deconvolution}
}

\pagestyle{fancy}
\fancyhf{}
\fancyhead[L]{\small Fabius--Thue--Morse convergence}
\fancyhead[R]{\small\nouppercase{\leftmark}}
\fancyfoot[C]{\thepage}
\renewcommand{\headrulewidth}{0.4pt}

\titleformat{\section}{\Large\bfseries}{\thesection}{0.7em}{}
\titleformat{\subsection}{\large\bfseries}{\thesubsection}{0.7em}{}
\titleformat{\subsubsection}{\normalsize\bfseries}{\thesubsubsection}{0.7em}{}
\setcounter{secnumdepth}{3}
\setcounter{tocdepth}{2}
\setlist[itemize]{leftmargin=2em,itemsep=0.25em,topsep=0.4em}
\setlist[enumerate]{leftmargin=2.3em,itemsep=0.25em,topsep=0.4em}
\setlength{\emergencystretch}{2em}

\newtheorem{theorem}{Theorem}[section]
\newtheorem{proposition}[theorem]{Proposition}
\newtheorem{lemma}[theorem]{Lemma}
\newtheorem{corollary}[theorem]{Corollary}
\newtheorem{conjecture}[theorem]{Conjecture}
\theoremstyle{definition}
\newtheorem{definition}[theorem]{Definition}
\newtheorem{algorithm}[theorem]{Algorithm}
\newtheorem{example}[theorem]{Example}
\theoremstyle{remark}
\newtheorem{remark}[theorem]{Remark}
\newtheorem{warning}[theorem]{Warning}

\newcommand{\R}{\mathbb R}
\newcommand{\C}{\mathbb C}
\newcommand{\N}{\mathbb N}
\newcommand{\Z}{\mathbb Z}
\newcommand{\Up}{\operatorname{up}}
\newcommand{\wt}{\operatorname{wt}_2}
\newcommand{\e}{\mathrm e}
\newcommand{\ii}{\mathrm i}
\newcommand{\dd}{\,\mathrm d}
\newcommand{\Var}{\operatorname{Var}}
\newcommand{\E}{\mathbb E}
\newcommand{\bigO}{\mathcal O}
\newcommand{\smallo}{o}
\newcommand{\supp}{\operatorname{supp}}
\newcommand{\dist}{\operatorname{dist}}
\newcommand{\repo}{\nolinkurl{Analysis/FabiusFunction}}
\newcommand{\fract}[1]{\left\{#1\right\}}
\newcommand{\Deltah}{\Delta_h}

\newenvironment{proofidea}{\par\smallskip\noindent\textit{Proof architecture.}\ }{\hfill$\square$\par\smallskip}
\newenvironment{boxedremark}{%
  \par\smallskip\noindent\begin{center}\begin{minipage}{0.94\textwidth}%
  \color{black}\hrule\vspace{0.6em}\small
}{%
  \vspace{0.6em}\hrule\end{minipage}\end{center}\smallskip
}

\lstdefinestyle{pseudo}{
  basicstyle=\ttfamily\small,
  backgroundcolor=\color{shadegray},
  frame=single,
  rulecolor=\color{rulegray},
  columns=fullflexible,
  keepspaces=true,
  showstringspaces=false,
  xleftmargin=0.7em,
  xrightmargin=0.7em,
  aboveskip=0.8em,
  belowskip=0.8em
}

\title{\bfseries Signed Thue--Morse Prefix Sums and the Fabius Function\\[0.35em]
\large Sharp convergence rates, lattice deconvolution, and exact grid drift}
\author{A quantitative supplement to the formal development\\
\small Vladimir Reshetnikov's \texttt{ProveIt} repository}
\date{Repository snapshot: 25 August 2026}

\begin{document}
\maketitle

\begin{abstract}
Let $\varepsilon_m=(-1)^{\wt(m)}$ be the signed Thue--Morse sequence and let
$\sigma_{k+1}(m)=\sum_{j=0}^{m}\sigma_k(j)$ be its inclusive iterated prefix sums.  The
formal development in \repo\ proves that a shifted and normalized order-$(n+1)$ prefix row
is exactly the coefficient row of a finite product of discrete uniform laws, and that its
associated histogram converges pointwise to Rvachev's compactly supported function $\Up$.
Sampling the same row at $\lfloor 2^n x\rfloor$ consequently converges to the bounded Fabius
function $F(x)$.  The existing argument is deliberately qualitative and supplies no rate.

This report makes the convergence quantitative.  At the natural cell centers, the normalized
prefix coefficients admit a uniform deconvolution expansion
\[
 H_n(\lambda)=\Up(\lambda)
 -\frac{3n+4}{72\,4^n}\Up''(\lambda)
 +\left(\frac{15n+16}{43200\,16^n}
       +\frac{(3n+4)^2}{10368\,16^n}\right)\Up^{(4)}(\lambda)
 +\bigO\!\left(\frac{n^3}{64^n}\right).
\]
It follows that the sharp uniform center-grid error is
$\frac13 n4^{-n}(1+o(1))$; the same rate holds after piecewise-linear interpolation.
The piecewise-constant histogram is less accurate because of cell quantization:
$\|\mathcal W_n-\Up\|_\infty=2^{-n}+\bigO(n4^{-n})$.
Most strikingly, the literal dyadic-floor approximation
\[
 A_n(x)=\frac{\sigma_{n+1}(\lfloor2^n x\rfloor)}{2^{\binom n2}}
\]
has a still larger, entirely geometric first-order error.  Uniformly on $[0,1]$,
\[
 A_n(x)=F(x)+\frac{n+2-2\fract{2^nx}}{2^{n+1}}F'(x)
          +\bigO\!\left(\frac{n^2}{4^n}\right),
\]
and therefore
\[
 \left\|\frac{2^{n+1}}n(A_n-F)-F'\right\|_\infty\longrightarrow0,
 \qquad
 \frac{2^n}{n}\|A_n-F\|_\infty\longrightarrow1.
\]
The factor $n$ comes from the exact inward displacement of the finite support, not from slow
weak convergence.  At $x=1/2$ the phenomenon has the closed form
$A_n(1/2)-1/2=(n+2)2^{-n}-2^{-\binom n2}$.

The bridge identities used as input are formalized in Lean.  The Fourier asymptotic expansion,
its sharp constants, and the accelerated finite-difference constructions developed here are
new analytic deductions and are not claimed to have been formalized in the repository.
\end{abstract}

\tableofcontents
\newpage

\section{Scope, conventions, and the answer in one table}
\label{sec:scope}

The source article \cite{reshetnikovExposition} isolates the connection between iterated
prefix sums and $\Up$, corrects two index/normalization defects in a naive formulation, and
proves pointwise convergence.  It also says explicitly that its monotonicity-and-continuity
argument gives no uniform rate.  The purpose of this supplement is to determine that rate,
including its leading constant, and to explain why several superficially similar readings of
the same integer row have different orders of accuracy.

We use the source article's normalization throughout:
\begin{itemize}
\item $F$ is the bounded distribution function, equal to $0$ on $(-\infty,0]$ and to $1$ on
      $[1,\infty)$;
\item $\Up$ is the even compactly supported density on $[-1,1]$, with
      $\Up(x-1)=F(x)$ for $0\le x\le1$;
\item the Fourier transform is
      $\widehat f(\xi)=\int_{\R}f(x)\e^{-2\pi\ii x\xi}\dd x$.
\end{itemize}
The sharp derivative theorem from the source gives
\begin{equation}
 \|F^{(r)}\|_\infty=\|\Up^{(r)}\|_\infty
 =2^{\binom{r+1}{2}}.                                      \label{eq:sharp-derivatives}
\end{equation}
In particular $\|F'\|_\infty=\|\Up'\|_\infty=2$ and
$\|\Up''\|_\infty=8$.

\begin{boxedremark}
\textbf{The three natural rates.}
Put $h_n=2^{-n}$.  The same normalized prefix row produces three different errors.
\begin{center}
\begin{tabularx}{0.95\textwidth}{@{}>{\raggedright\arraybackslash}p{0.23\textwidth}
>{\raggedright\arraybackslash}p{0.30\textwidth}X@{}}
\toprule
Reading of the row & Sharp uniform error & Dominant mechanism\\
\midrule
Natural center values, or their linear interpolant
& $\displaystyle \frac13 n h_n^2(1+o(1))$
& Deconvolution of the omitted mean-zero Bernoulli tail; its variance is of order $nh_n^2$.\\[0.35em]
Piecewise-constant histogram $\mathcal W_n$
& $\displaystyle h_n+\bigO(nh_n^2)$
& Horizontal quantization by half a cell, amplified by the optimal slope $2$.\\[0.35em]
Dyadic-floor Fabius sample $A_n$
& $\displaystyle n h_n(1+o(1))$
& The finite support is shorter than $[-1,1]$ by $(n+2)h_n/2$ at each end, producing an inward drift.\\
\bottomrule
\end{tabularx}
\end{center}
Thus the pointwise floor formula is not the most accurate reconstruction of the function from
the prefix row.  It is worse than the histogram by a factor asymptotic to $n$, while a
centered linear reconstruction is better than the histogram by a factor asymptotic to
$3\cdot2^n/n$.
\end{boxedremark}

\subsection{What is formalized and what is new here}

The exact prefix convolution law, Prouhet cancellation, zero runs, generating-series identity,
coefficient bridge, corrected cell centers, and pointwise limit are developed in
\path{ThueMorsePrefix.lean} and \path{ThueMorseApproximation.lean}; the finite atomic and step
approximants are developed in \path{EarlyApproximants.lean},
\path{EarlyMeasureBridge.lean}, \path{StepMeasureBridge.lean}, and
\path{StepApproximationLimit.lean} \cite{proveit}.  These are the algebraic and
measure-theoretic inputs used below.

The new ingredient is a quantitative inversion of the finite cosine product on its natural
Nyquist interval.  It yields a complete asymptotic expansion in even derivatives of $\Up$.
The first two nonzero terms are enough to determine all sharp rates stated in the abstract.
No part of that new expansion is represented here as already machine-checked.

\section{The exact discrete bridge}
\label{sec:bridge}

\subsection{Signed Thue--Morse prefixes}

Let
\begin{equation}
 \varepsilon_m=(-1)^{\wt(m)},\qquad m\ge0,                 \label{eq:tm-sign}
\end{equation}
where $\wt(m)$ is the number of ones in the binary expansion of $m$.  Define inclusive
iterated prefix sums by
\begin{equation}
 \sigma_0(m)=\varepsilon_m,
 \qquad
 \sigma_{k+1}(m)=\sum_{j=0}^{m}\sigma_k(j).                \label{eq:prefix-def}
\end{equation}
The word \emph{inclusive} matters: one has
\begin{equation}
 \sigma_{k+1}(m+1)-\sigma_{k+1}(m)=\sigma_k(m+1),          \label{eq:forward-difference}
\end{equation}
with the lower-order value at the new index.

Repeated summation is convolution with the discrete truncated-power kernel.  For $k\ge1$,
\begin{equation}
 \sigma_k(m)=\sum_{r=0}^{m}\binom{m-r+k-1}{k-1}\varepsilon_r. \label{eq:prefix-convolution}
\end{equation}
The signed generating series is
\begin{equation}
 \sum_{m\ge0}\varepsilon_m z^m=\prod_{j\ge0}(1-z^{2^j}),
 \qquad
 \sum_{m\ge0}\sigma_k(m)z^m
 =\frac{\prod_{j\ge0}(1-z^{2^j})}{(1-z)^k}.               \label{eq:prefix-generating}
\end{equation}
All identities in this paragraph are formal power-series identities over $\Z$.

\subsection{The finite polynomial and the one-order shift}

For $n\ge0$ put
\begin{equation}
 P_n(z)=\prod_{i=1}^{n}(1+z+\cdots+z^{2^i-1}),
 \qquad
 g_n=\deg P_n=2^{n+1}-n-2.                                \label{eq:Pn}
\end{equation}
The cancellation of the first $n+1$ factors in \eqref{eq:prefix-generating} gives the exact
coefficient identity
\begin{equation}
 [z^m]P_n(z)=\sigma_{n+1}(m),
 \qquad 0\le m<2^{n+1}.                                   \label{eq:coefficient-bridge}
\end{equation}
This is the source of the order shift: the polynomial indexed by $n$ corresponds to the
$(n+1)$st, not the $n$th, prefix row.

Every factor of $P_n$ is palindromic, hence
\begin{equation}
 [z^m]P_n=[z^{g_n-m}]P_n,
 \qquad
 P_n(1)=2^{\binom{n+1}{2}}.                               \label{eq:Pn-symmetry-mass}
\end{equation}
The coefficients are nonnegative integers.  Thus, after normalization, they form a symmetric
probability mass function.

\subsection{Natural centers, atomic law, and histogram}

Set
\begin{equation}
 h_n=2^{-n},
 \qquad
 \lambda_{n,m}=\frac{2m-g_n}{2^{n+1}}
 =h_n\left(m-\frac{g_n}{2}\right).                        \label{eq:center}
\end{equation}
The points $\lambda_{n,m}$ lie on a translate of the lattice $h_n\Z$.  Define
\begin{equation}
 p_{n,m}=2^{-\binom{n+1}{2}}[z^m]P_n,
 \qquad
 \mu_n=\sum_{m\in\Z}p_{n,m}\,\delta_{\lambda_{n,m}},    \label{eq:atomic-law}
\end{equation}
where coefficients outside $0\le m\le g_n$ are understood to be zero.  Then $\mu_n$ is a
probability measure.  Its mass density relative to the lattice spacing is
\begin{equation}
 H_n(\lambda_{n,m})=\frac{p_{n,m}}{h_n}
 =\frac{[z^m]P_n}{2^{\binom n2}}.                          \label{eq:Hn}
\end{equation}
By \eqref{eq:coefficient-bridge}, this is exactly
\begin{equation}
 H_n(\lambda_{n,m})
 =\frac{\sigma_{n+1}(m)}{2^{\binom n2}},
 \qquad 0\le m<2^{n+1}.                                  \label{eq:Hn-prefix}
\end{equation}

Let $I_{n,m}$ be the cell of width $h_n$ centered at $\lambda_{n,m}$.  Spreading each atom of
$\mu_n$ uniformly over its cell gives the step density
\begin{equation}
 \mathcal W_n(x)=\sum_m H_n(\lambda_{n,m})\mathbf 1_{I_{n,m}}(x), \label{eq:Wn}
\end{equation}
with half weight from each adjacent cell at a shared endpoint.  In the interior of a cell,
$\mathcal W_n$ is simply the corresponding normalized prefix coefficient.  The formal
repository proves
\begin{equation}
 \mathcal W_n(x)\longrightarrow\Up(x)                     \label{eq:qualitative-W}
\end{equation}
pointwise.  The rest of this report quantifies \eqref{eq:qualitative-W} and the moving-center
sample derived from it.

\section{The hidden Bernoulli tail}
\label{sec:tail}

\subsection{Rademacher factorization}

Let $R_{\ell,r}$ be independent Rademacher variables, each taking $\pm1$ with probability
$1/2$.  The binary factorization of every discrete uniform variable in \eqref{eq:Pn} gives
an equivalent representation of \eqref{eq:atomic-law}:
\begin{equation}
 X_n=\sum_{\ell=1}^{n}\sum_{r=1}^{\ell}
       2^{-\ell-1}R_{\ell,r},
 \qquad \mathcal L(X_n)=\mu_n.                            \label{eq:Xn}
\end{equation}
Consequently
\begin{equation}
 \widehat{\mu_n}(\xi)
 =\prod_{\ell=1}^{n}
  \left(\cos\frac{\pi\xi}{2^\ell}\right)^\ell.          \label{eq:finite-cosine}
\end{equation}
The complete series
\begin{equation}
 Y=\sum_{\ell=1}^{\infty}\sum_{r=1}^{\ell}
       2^{-\ell-1}R_{\ell,r}                              \label{eq:Y}
\end{equation}
converges absolutely, and its density is $\Up$.  Equivalently,
\begin{equation}
 \widehat\Up(\xi)
 =\prod_{\ell=1}^{\infty}
  \left(\cos\frac{\pi\xi}{2^\ell}\right)^\ell.          \label{eq:infinite-cosine}
\end{equation}
This cosine representation is equivalent to the more familiar infinite sinc product.

Write
\begin{equation}
 T_n=Y-X_n
 =\sum_{\ell>n}\sum_{r=1}^{\ell}2^{-\ell-1}R_{\ell,r}.    \label{eq:Tn}
\end{equation}
Then $X_n$ and $T_n$ are independent and
\begin{equation}
 \Up(x)\dd x=\mu_n*\mathcal L(T_n).                        \label{eq:tail-convolution}
\end{equation}
Thus $\mu_n$ is obtained by deleting a small, symmetric smoothing tail from the target law.
The center-grid asymptotics are an exact lattice version of deconvolving that tail.

\subsection{Exact support radius and variance}

\begin{proposition}[Geometry and second moment of the omitted tail]
\label{prop:tail-moments}
The tail $T_n$ has mean zero, support contained in $[-r_n,r_n]$, and variance $v_n$, where
\begin{align}
 r_n&=\sum_{\ell>n}\ell2^{-\ell-1}
     =\frac{n+2}{2^{n+1}},                                  \label{eq:rn}\\
 v_n&=\sum_{\ell>n}\ell2^{-2\ell-2}
     =\frac{3n+4}{36\,4^n}.                                 \label{eq:vn}
\end{align}
Both support endpoints occur.
\end{proposition}

\begin{proof}
The mean vanishes by symmetry.  The maximal possible absolute value is obtained by taking
all signs equal, so it is the first displayed sum.  The geometric-tail identity
\[
 \sum_{\ell=n+1}^{\infty}\ell q^\ell
 =q^{n+1}\frac{(n+1)-nq}{(1-q)^2}
\]
with $q=1/2$ gives \eqref{eq:rn}.  Independence and $\Var(R_{\ell,r})=1$ give the variance
sum, and the same identity with $q=1/4$, followed by the extra factor $1/4$, gives
\eqref{eq:vn}.
\end{proof}

The support of $\mu_n$ is therefore
\begin{equation}
 \supp\mu_n=[-1+r_n,1-r_n]\cap\bigl(h_n(\Z-g_n/2)\bigr),  \label{eq:support-mun}
\end{equation}
because $g_n/2^{n+1}=1-r_n$.  Two scales should be distinguished:
\begin{equation}
 r_n\sim\frac n2h_n,
 \qquad
 \sqrt{v_n}\sim\sqrt{\frac n{12}}\,h_n.                  \label{eq:two-tail-scales}
\end{equation}
The support radius is a worst-case all-signs quantity, while the standard deviation is the
typical scale.  The floor-indexed approximation will inherit the former; the correctly
centered deconvolution error will inherit the variance itself.

\subsection{Why the leading center correction is second order}

If a smooth density $f$ is convolved with a small mean-zero law of variance $v$, then formally
\[
 f*\nu=f+\frac v2f''+\cdots.
\]
Solving backward gives $f-\frac v2f''+\cdots$.  Since \eqref{eq:tail-convolution} smooths the
atomic row by the tail law, one should expect
\begin{equation}
 H_n(\lambda)-\Up(\lambda)
 \approx-\frac{v_n}{2}\Up''(\lambda)
 =-\frac{3n+4}{72\,4^n}\Up''(\lambda).                    \label{eq:heuristic-leading}
\end{equation}
The sign is that of deconvolution: in a convex region the unsmoothed profile lies below the
smoothed one.  The next sections make \eqref{eq:heuristic-leading} uniform and rigorous and
compute the fourth-derivative term.

\section{Fourier inversion on the moving lattice}
\label{sec:fourier}

The main technical point is that one must invert the finite atomic law on exactly one period
of its lattice Fourier series.  Ordinary Fourier inversion of a density cannot be applied
directly to $\mu_n$.

\subsection{A shifted-lattice inversion formula}

\begin{lemma}[Nyquist inversion for a shifted lattice]
\label{lem:lattice-inversion}
Let $\mu=\sum_{k\in\Z}p_k\delta_{a+hk}$ be a finite measure on a translate of $h\Z$, and let
$\widehat\mu(\xi)=\sum_kp_k\e^{-2\pi\ii(a+hk)\xi}$.  Then, for every $m\in\Z$,
\begin{equation}
 \frac{p_m}{h}
 =\int_{-1/(2h)}^{1/(2h)}
   \widehat\mu(\xi)\e^{2\pi\ii(a+hm)\xi}\dd\xi.          \label{eq:lattice-inversion}
\end{equation}
\end{lemma}

\begin{proof}
Insert the finite sum defining $\widehat\mu$ and integrate term by term.  The $k$th term is
\[
 p_k\int_{-1/(2h)}^{1/(2h)}\e^{2\pi\ii h(m-k)\xi}\dd\xi,
\]
which is $p_m/h$ when $k=m$ and zero otherwise.  The shift $a$ cancels from the phase.
\end{proof}

For $\mu_n$, the interval in \eqref{eq:lattice-inversion} is
\begin{equation}
 B_n=[-2^{n-1},2^{n-1}].                                   \label{eq:Bn}
\end{equation}
Combining \eqref{eq:Hn} and \eqref{eq:lattice-inversion} gives
\begin{equation}
 H_n(\lambda_{n,m})
 =\int_{B_n}\widehat{\mu_n}(\xi)
   \e^{2\pi\ii\lambda_{n,m}\xi}\dd\xi.                  \label{eq:Hn-inversion}
\end{equation}

\subsection{The exact deconvolution multiplier}

On $B_n$, all factors in the omitted cosine tail are positive: for $\ell>n$,
$|\pi\xi/2^\ell|\le\pi/4$.  Hence one may divide \eqref{eq:infinite-cosine} by
\eqref{eq:finite-cosine} and write
\begin{equation}
 \widehat{\mu_n}(\xi)=\widehat\Up(\xi)R_n(\xi),
 \qquad
 R_n(\xi)=\prod_{\ell>n}
  \left(\sec\frac{\pi\xi}{2^\ell}\right)^\ell,
 \qquad \xi\in B_n.                                      \label{eq:Rn}
\end{equation}
The multiplier $R_n$ is the reciprocal characteristic function of the omitted tail on the
frequency range where lattice inversion needs it.

There is a useful global bound.  Put $t=\pi\xi/2^n$, so $|t|\le\pi/2$.  Reindexing
$\ell=n+r$ and using Vi\`ete's product
$\prod_{r\ge1}\cos(t/2^r)=\sin t/t$ gives
\begin{equation}
 R_n(\xi)
 =\left(\frac{t}{\sin t}\right)^n
   \prod_{r=1}^{\infty}\sec(t/2^r)^r.                     \label{eq:Rn-viete}
\end{equation}
The second product is bounded uniformly for $|t|\le\pi/2$, while
$t/\sin t\le\pi/2$.  Thus there is an absolute $C_*$ such that
\begin{equation}
 1\le R_n(\xi)\le C_*\left(\frac\pi2\right)^n,
 \qquad \xi\in B_n.                                      \label{eq:Rn-global-bound}
\end{equation}
This exponential growth is harmless because $\widehat\Up$ decays faster than every power.
Indeed, compact support and smoothness give, for each $M\ge1$,
\begin{equation}
 |\widehat\Up(\xi)|
 \le\frac{\|\Up^{(M)}\|_1}{(2\pi|\xi|)^M}
 \le\frac{2^{1+\binom{M+1}{2}}}{(2\pi|\xi|)^M}
 \qquad(\xi\ne0),                                        \label{eq:fourier-rapid}
\end{equation}
where \eqref{eq:sharp-derivatives} was used in the second inequality.

\section{Sharp center-grid asymptotics}
\label{sec:center-rate}

\subsection{Expansion of the reciprocal tail product}

For $|u|<\pi/2$,
\begin{equation}
 \log\sec u=\frac{u^2}{2}+\frac{u^4}{12}+\bigO(u^6).       \label{eq:logsec}
\end{equation}
The two geometric tails needed below are
\begin{align}
 \frac12\sum_{\ell>n}\ell4^{-\ell}
 &=\frac{3n+4}{18\,4^n}=:\alpha_n,                        \label{eq:alpha-n}\\
 \frac1{12}\sum_{\ell>n}\ell16^{-\ell}
 &=\frac{15n+16}{2700\,16^n}=: \beta_n.                  \label{eq:beta-n}
\end{align}
Therefore, on a low-frequency range growing sufficiently slowly with $n$,
\begin{equation}
 \log R_n(\xi)
 =\alpha_n(\pi\xi)^2+\beta_n(\pi\xi)^4
  +\bigO\!\left(\frac{n}{64^n}|\xi|^6\right),            \label{eq:log-Rn}
\end{equation}
and exponentiation gives
\begin{equation}
 R_n(\xi)=1+\alpha_n(\pi\xi)^2
 +\left(\beta_n+\frac{\alpha_n^2}{2}\right)(\pi\xi)^4
 +\bigO\!\left(\frac{n^3}{64^n}|\xi|^6\right).          \label{eq:Rn-expansion}
\end{equation}

\subsection{The uniform expansion}

\begin{theorem}[Two-term lattice deconvolution expansion]
\label{thm:center-expansion}
Uniformly for all lattice centers $\lambda=\lambda_{n,m}$,
\begin{align}
 H_n(\lambda)
={}&\Up(\lambda)
 -\frac{3n+4}{72\,4^n}\Up''(\lambda) \notag\\
 &+\left(
   \frac{15n+16}{43200\,16^n}
  +\frac{(3n+4)^2}{10368\,16^n}
  \right)\Up^{(4)}(\lambda)
 +\bigO\!\left(\frac{n^3}{64^n}\right).                 \label{eq:center-expansion}
\end{align}
The implicit constant is absolute and effective.  In particular,
\begin{equation}
 H_n(\lambda)-\Up(\lambda)
 =-\frac{3n+4}{72\,4^n}\Up''(\lambda)
  +\bigO\!\left(\frac{n^2}{16^n}\right)                 \label{eq:center-leading}
\end{equation}
uniformly on the moving lattice.
\end{theorem}

\begin{proof}
Use \eqref{eq:Hn-inversion} and \eqref{eq:Rn} and split the integral at
$|\xi|=2^{n/4}$.  On the inner interval, \eqref{eq:Rn-expansion} is uniform.  Since
$\int_{\R}|\xi|^6|\widehat\Up(\xi)|\dd\xi<\infty$, its integrated remainder is
$\bigO(n^3 64^{-n})$ uniformly in the phase $\lambda$.

On the complementary part of $B_n$, combine \eqref{eq:Rn-global-bound} with
\eqref{eq:fourier-rapid}.  Taking, for example, $M=32$ makes the resulting integral
$o(64^{-n})$.  The tails of the three polynomial Fourier integrals outside $B_n$ are even
smaller.  We may therefore replace the truncated integrals of the displayed polynomial terms
by integrals over all of $\R$.

With the chosen Fourier convention,
\begin{equation*}
 \int_{\R}(\pi\xi)^2\widehat\Up(\xi)\e^{2\pi\ii\lambda\xi}\dd\xi
 =-\frac14\Up''(\lambda),
 \qquad
 \int_{\R}(\pi\xi)^4\widehat\Up(\xi)\e^{2\pi\ii\lambda\xi}\dd\xi
 =\frac1{16}\Up^{(4)}(\lambda).
\end{equation*}
Substituting \eqref{eq:alpha-n} and \eqref{eq:beta-n} gives
\eqref{eq:center-expansion}.  Omitting the fourth-derivative term and using
\eqref{eq:sharp-derivatives} gives \eqref{eq:center-leading}.
\end{proof}

\begin{remark}[Variance form]
Because $\alpha_n=2v_n$, the leading term is exactly
\begin{equation}
 H_n=\Up-\frac{v_n}{2}\Up''+\cdots,                        \label{eq:variance-form}
\end{equation}
which confirms the deconvolution heuristic \eqref{eq:heuristic-leading}.
\end{remark}

\subsection{Sharp norm constant}

\begin{corollary}[Sharp center-grid rate]
\label{cor:center-sharp}
One has the uniform profile limit
\begin{equation}
 \sup_{m\in\Z}
 \left|
  \frac{4^n}{n}\bigl(H_n(\lambda_{n,m})-\Up(\lambda_{n,m})\bigr)
  +\frac1{24}\Up''(\lambda_{n,m})
 \right|\longrightarrow0.                                \label{eq:center-profile-limit}
\end{equation}
Consequently,
\begin{equation}
 \lim_{n\to\infty}\frac{4^n}{n}
 \sup_{m\in\Z}|H_n(\lambda_{n,m})-\Up(\lambda_{n,m})|
 =\frac13.                                                 \label{eq:center-sharp-norm}
\end{equation}
\end{corollary}

\begin{proof}
The coefficient $(3n+4)/(72n)$ tends to $1/24$, while the remainder in
\eqref{eq:center-leading}, after multiplication by $4^n/n$, tends uniformly to zero.  The
lattices become dense, and $\Up''$ is continuous with exact supremum $8$.  This supremum
is attained; for example, the attained-point form of the sharp derivative theorem gives
$\Up''(-3/4)=F''(1/4)=8$.  Hence the norm limit is $8/24=1/3$.
\end{proof}

\subsection{A continuous reconstruction with the same rate}

Let $L_n$ be the piecewise-linear interpolant through the points
$(\lambda_{n,m},H_n(\lambda_{n,m}))$, with zero values continued outside the finite support.
The ordinary interpolation error for a $C^2$ function is at most
$h_n^2\|\Up''\|_\infty/8=h_n^2$.  This is smaller by a factor $1/n$ than the leading
center bias.

\begin{corollary}[Sharp continuous linear reconstruction]
\label{cor:linear-rate}
Uniformly on $\R$,
\begin{equation}
 \frac{4^n}{n}(L_n-\Up)\longrightarrow-\frac1{24}\Up'',   \label{eq:linear-profile}
\end{equation}
in the supremum norm.  In particular,
\begin{equation}
 \|L_n-\Up\|_\infty=\frac n{3\,4^n}(1+o(1)).              \label{eq:linear-sharp}
\end{equation}
\end{corollary}

Thus iterated prefix sums do not merely converge to $\Up$ after a qualitative passage through
histograms.  Read at their actual centers and interpolated linearly, they give a continuous
approximation with a fully determined, curvature-shaped leading error.

\subsection{The all-orders pattern}

The same proof is not tied to fourth order.  Write
\begin{equation}
 \log\sec z=\sum_{r\ge1}c_rz^{2r}
 \quad(c_1=1/2,\ c_2=1/12,\ c_3=1/45,\ldots)              \label{eq:logsec-full}
\end{equation}
and put
\begin{equation}
 s_{r,n}=c_r\sum_{\ell>n}\ell2^{-2r\ell}.                 \label{eq:srn}
\end{equation}
For a fixed $M$, expand
\begin{equation}
 \exp\left(\sum_{r=1}^{M-1}s_{r,n}z^{2r}\right)
 =\sum_{j=0}^{M-1}b_{j,n}z^{2j}+\bigO(z^{2M}).             \label{eq:bjn}
\end{equation}
Then the lattice inversion argument gives
\begin{equation}
 H_n(\lambda)
 =\sum_{j=0}^{M-1}\frac{(-1)^j b_{j,n}}{4^j}
    \Up^{(2j)}(\lambda)
 +\bigO_M\!\left(n^M4^{-Mn}\right),                      \label{eq:all-orders}
\end{equation}
uniformly on the center lattice.  Each $b_{j,n}$ is an explicit polynomial in the geometric
tails $s_{r,n}$.  Thus the connection admits a full asymptotic differential expansion, not
just a big-$O$ convergence estimate.

\section{The histogram: a sharp first-order quantization error}
\label{sec:histogram-rate}

The center values are exceptionally accurate, but $\mathcal W_n$ holds each value constant
across a cell.  That reconstruction injects a larger first-order horizontal error.

For a point $x$ in the interior of a cell, let $c_n(x)$ denote that cell's center and define
the sawtooth phase
\begin{equation}
 \rho_n(x)=\frac{c_n(x)-x}{h_n},
 \qquad -\frac12<\rho_n(x)<\frac12.                       \label{eq:rho}
\end{equation}
Then \eqref{eq:center-leading} and Taylor's theorem give, uniformly away from the shared cell
endpoints,
\begin{equation}
 \mathcal W_n(x)-\Up(x)
 =h_n\rho_n(x)\Up'(x)+\bigO(nh_n^2).                      \label{eq:hist-local}
\end{equation}
At a shared endpoint, the half-weight convention averages the two adjacent center values and
cancels the first-order odd displacement, leaving only $\bigO(nh_n^2)$ there.

\begin{theorem}[Sharp histogram rate]
\label{thm:hist-sharp}
The step approximants satisfy
\begin{equation}
 \|\mathcal W_n-\Up\|_\infty
 =2^{-n}+\bigO\!\left(\frac n{4^n}\right).                \label{eq:hist-sharp}
\end{equation}
Equivalently,
\begin{equation}
 \lim_{n\to\infty}2^n\|\mathcal W_n-\Up\|_\infty=1.      \label{eq:hist-limit-constant}
\end{equation}
\end{theorem}

\begin{proof}
For an interior point, combine \eqref{eq:center-leading} with the optimal Lipschitz bound
$\|\Up'\|_\infty=2$:
\[
 |\mathcal W_n(x)-\Up(x)|
 \le \bigO(nh_n^2)+2|c_n(x)-x|
 \le h_n+\bigO(nh_n^2).
\]
At a shared endpoint the averaged value obeys the same bound.  There remains only the thin
uncovered strip between the support of $\mathcal W_n$ and an endpoint of $[-1,1]$.
Its width is at most $r_n+h_n/2=\bigO(nh_n)$.  Since $\Up$ and all of its
derivatives vanish at $\pm1$, Taylor's theorem with the sharp bounded third derivative gives
\[
 |\Up(\pm1\mp s)|\le \frac{\|\Up^{(3)}\|_\infty}{6}s^3
 =\bigO(n^3h_n^3)=\smallo(nh_n^2)
 \qquad(0\le s\le r_n+h_n/2).
\]
Outside $[-1,1]$ both functions vanish.  This completes the upper estimate globally.

For the lower estimate, choose cells whose centers tend to $-1/2$, where
$\Up'(-1/2)=2$, and approach one edge of each chosen cell from its interior.  The center error
is $\bigO(nh_n^2)$, while the change of $\Up$ across half a cell is
$h_n+\bigO(h_n^2)$.  Taking the supremum gives
$\|\mathcal W_n-\Up\|_\infty\ge h_n-\bigO(nh_n^2)$.
\end{proof}

\begin{remark}
The pointwise proof in the source article could not reveal \eqref{eq:hist-sharp}: weak
convergence plus unimodality controls values through fixed comparison points but does not
track a moving point at distance $h_n/2$ from a cell center.  The rate is a local geometric
phenomenon, and the sharp constant comes from the sharp derivative bound.
\end{remark}

\section{The floor-prefix approximation to the bounded Fabius function}
\label{sec:floor-rate}

\subsection{Definition and exact moving-center displacement}

The source theorem obtains $F$ from the prefix row by evaluating its $\Up$ approximation at
half the argument.  Define
\begin{equation}
 A_n(x)=\mathcal J_n(x/2)
 =\frac{\sigma_{n+1}(\lfloor2^nx\rfloor)}{2^{\binom n2}},
 \qquad 0\le x\le1.                                      \label{eq:An}
\end{equation}
Let
\begin{equation}
 \theta_n(x)=\fract{2^nx}=2^nx-\lfloor2^nx\rfloor.        \label{eq:theta}
\end{equation}
The normalized prefix value in \eqref{eq:An} is not located at $x-1$.  Its actual center is
\begin{align}
 \lambda_n(x)
 &=\lambda_{n,\lfloor2^nx\rfloor} \notag\\
 &=x-1+\delta_n(x),                                        \label{eq:lambda-drift}\\
 \delta_n(x)
 &=\frac{n+2-2\theta_n(x)}{2^{n+1}}.                       \label{eq:delta}
\end{align}
This identity, not merely the fact that $\lambda_n(x)\to x-1$, determines the dominant error.
It also gives
\begin{equation}
 \frac{n}{2^{n+1}}<\delta_n(x)\le\frac{n+2}{2^{n+1}}.     \label{eq:delta-range}
\end{equation}
In terms of the omitted-tail support radius \eqref{eq:rn},
\begin{equation}
 \delta_n(x)=r_n-\theta_n(x)h_n.                           \label{eq:delta-rn}
\end{equation}
Thus the index $\lfloor2^nx\rfloor$ compensates for the ordinary fractional cell position but
not for the shortening of the finite support.  The residual is roughly $r_n\sim nh_n/2$,
which is much larger than either the cell width error of the histogram or the variance-scale
center bias.

\subsection{First- and second-order expansions}

\begin{theorem}[Uniform floor-prefix expansion]
\label{thm:floor-expansion}
Uniformly for $x\in[0,1]$,
\begin{equation}
 A_n(x)
 =F(x)+\delta_n(x)F'(x)
  +\bigO\!\left(\frac{n^2}{4^n}\right).                  \label{eq:floor-first}
\end{equation}
More precisely,
\begin{align}
 A_n(x)
={}&F(x)+\delta_n(x)F'(x) \notag\\
 &+\left(
    \frac{\delta_n(x)^2}{2}
    -\frac{3n+4}{72\,4^n}
   \right)F''(x)
 +\bigO\!\left(\frac{n^3}{8^n}\right).                  \label{eq:floor-second}
\end{align}
All remainder constants are independent of $x$ and $n$.
\end{theorem}

\begin{proof}
By the exact coefficient identification,
$A_n(x)=H_n(\lambda_n(x))$.  Apply \eqref{eq:center-leading} at
$\lambda_n(x)=x-1+\delta_n(x)$, use $\Up(x-1)=F(x)$ and the same identity for derivatives,
and Taylor-expand around $x-1$.  Since $\delta_n=\bigO(n2^{-n})$, the first omitted Taylor
term in \eqref{eq:floor-first} is $\bigO(n^2 4^{-n})$.

For \eqref{eq:floor-second}, retain the quadratic Taylor term and the leading center-bias
term.  The remainders are
$\bigO(\delta_n^3)$, $\bigO(n4^{-n}\delta_n)$, and
$\bigO(n^2 16^{-n})$, all absorbed by $\bigO(n^3 8^{-n})$.
\end{proof}

\begin{corollary}[Uniform profile and sharp norm rate]
\label{cor:floor-sharp}
The entire rescaled error profile converges uniformly:
\begin{equation}
 \left\|\frac{2^{n+1}}n(A_n-F)-F'\right\|_\infty
 \longrightarrow0.                                       \label{eq:floor-profile}
\end{equation}
Consequently,
\begin{equation}
 \lim_{n\to\infty}\frac{2^n}{n}\|A_n-F\|_\infty=1.       \label{eq:floor-sharp}
\end{equation}
For each fixed $x$,
\begin{equation}
 \lim_{n\to\infty}\frac{2^{n+1}}n\bigl(A_n(x)-F(x)\bigr)
 =F'(x).                                                   \label{eq:floor-pointwise-profile}
\end{equation}
\end{corollary}

\begin{proof}
By \eqref{eq:delta},
\[
 \frac{2^{n+1}}n\delta_n(x)
 =1+\frac{2-2\theta_n(x)}n\longrightarrow1
\]
uniformly in $x$.  Multiplying the remainder in \eqref{eq:floor-first} by
$2^{n+1}/n$ gives $\bigO(n2^{-n})$.  This proves \eqref{eq:floor-profile}.  Taking norms and
using $\|F'\|_\infty=2$ gives \eqref{eq:floor-sharp}.
\end{proof}

\subsection{The equivalent statement for the \texorpdfstring{$\Up$}{Up} sample}

The source notation uses
\begin{equation}
 \mathcal J_n(u)
 =\frac{\sigma_{n+1}(\lfloor2^{n+1}u\rfloor)}{2^{\binom n2}}.
\end{equation}
For $0\le u<1$ put $\vartheta_n(u)=\fract{2^{n+1}u}$ and
\begin{equation}
 d_n(u)=\frac{n+2-2\vartheta_n(u)}{2^{n+1}}.
\end{equation}
Then
\begin{equation}
 \mathcal J_n(u)
 =\Up(2u-1)+d_n(u)\Up'(2u-1)
  +\bigO\!\left(\frac{n^2}{4^n}\right).                  \label{eq:J-expansion}
\end{equation}
At $u=1$, the prefix coefficient identity is outside its range and the exact zero-run value is
$-2^{-\binom n2}$; since $\Up(1)=\Up'(1)=0$, this exceptional endpoint is much smaller than
the uniform principal scale and does not change the norm asymptotic.

\section{An exact midpoint identity}
\label{sec:midpoint}

The sharp constant in \eqref{eq:floor-sharp} can be seen without asymptotic analysis at the
single point where $F'$ attains its maximum.

\begin{theorem}[Exact floor-prefix error at $x=1/2$]
\label{thm:midpoint-exact}
For every $n\ge2$,
\begin{equation}
 A_n\!\left(\frac12\right)
 =\frac12+\frac{n+2}{2^n}-2^{-\binom n2}.                 \label{eq:midpoint-exact}
\end{equation}
Equivalently,
\begin{equation}
 \frac{2^n}{n}
 \left(A_n\!\left(\frac12\right)-\frac12\right)
 =1+\frac2n-\frac{2^{n-\binom n2}}n.                     \label{eq:midpoint-scaled}
\end{equation}
\end{theorem}

\begin{proof}
Put $M=2^{n-1}$ and write $Q=P_{n-1}=\sum_{r=0}^{d}q_rz^r$, where
$d=g_{n-1}=2^n-n-1$, extending $q_r$ by zero outside $0\le r\le d$.  Since $M<2^n$,
multiplication by the last factor of $P_n$ gives
\begin{equation}
 [z^M]P_n=\sum_{r=0}^{M}q_r=:C.                            \label{eq:midpoint-C}
\end{equation}
The polynomial $Q$ is palindromic and has total mass
\begin{equation}
 T=Q(1)=2^{\binom n2}.                                     \label{eq:midpoint-T}
\end{equation}
By symmetry,
\begin{equation}
 2C-T=\sum_{r=d-M}^{M}q_r.                                 \label{eq:midpoint-band}
\end{equation}
The band in \eqref{eq:midpoint-band} has $n+2$ coefficients.  To evaluate them, write
$Q=P_{n-2}(1+\cdots+z^{2^{n-1}-1})$.  Put
$T_0=P_{n-2}(1)=2^{\binom{n-1}{2}}$.  The $n$ interior coefficients of the band are all
$T_0$, because the convolution window contains all coefficients of $P_{n-2}$.  Each of the
two boundary coefficients is $T_0-1$: its window omits only the monic leading coefficient of
$P_{n-2}$, and palindromicity gives the same value at the opposite boundary.  Hence
\begin{equation}
 2C-T=(n+2)T_0-2.                                         \label{eq:midpoint-band-value}
\end{equation}
Since $T/T_0=2^{n-1}$,
\[
 \frac CT=\frac12+\frac{n+2}{2^n}-\frac1T.
\]
Finally, \eqref{eq:coefficient-bridge} and \eqref{eq:An} identify $C/T$ with
$A_n(1/2)$, proving \eqref{eq:midpoint-exact}.
\end{proof}

\begin{table}[htbp]
\centering
\caption{Exact midpoint error and its sharp-rate normalization.}
\label{tab:midpoint}
\begin{tabular}{@{}rcc@{}}
\toprule
$n$ & $A_n(1/2)-1/2$ & $2^n\bigl(A_n(1/2)-1/2\bigr)/n$\\
\midrule
2  & $0.5$                 & $1.000000000000$\\
3  & $0.5$                 & $1.333333333333$\\
4  & $0.359375$            & $1.437500000000$\\
5  & $0.2177734375$        & $1.393750000000$\\
6  & $0.124969482421875$   & $1.333007812500$\\
8  & $0.0390624962747097$  & $1.249999880791$\\
10 & $0.0117187499999716$  & $1.199999999997$\\
12 & $0.0034179687500000$  & $1.166666666667$\\
16 & $0.0002746582031250$  & $1.125000000000$\\
20 & $0.00002098083496094$ & $1.100000000000$\\
\bottomrule
\end{tabular}
\end{table}

The exact identity also explains an otherwise puzzling numerical observation: after a few
levels, the error is almost indistinguishable from the tangent displacement
$2\delta_n(1/2)=(n+2)2^{-n}$.  The remaining defect is the super-exponentially small term
$2^{-\binom n2}$.

\subsection{Correct centering can be exact at the same point}

The large error in \eqref{eq:midpoint-exact} is not an unavoidable defect of the prefix row.
It is caused by reading the wrong index as the point $-1/2$ on the $\Up$ axis.  When $n$ is
even, $-1/2$ itself is a lattice center, with index
\begin{equation}
 m_n^*=2^{n-1}-\frac{n+2}{2}.                              \label{eq:midpoint-correct-index}
\end{equation}
For this index, $m_n^*=(g_{n-1}-1)/2$.  Since $P_{n-1}$ is palindromic of odd degree, exactly
half of its total coefficient mass lies at indices at most $m_n^*$.  The last-factor prefix
identity therefore gives
\begin{equation}
 H_n(-1/2)=\frac12=\Up(-1/2) \qquad(n\ \text{even}).       \label{eq:midpoint-centered-exact}
\end{equation}
The floor rule instead uses $M=2^{n-1}$, which is $(n+2)/2$ whole cells to the right of
$m_n^*$.  This is the discrete form of the support drift.

\section{Bias removal and accelerated prefix reconstructions}
\label{sec:acceleration}

The expansion \eqref{eq:center-expansion} can be used constructively.  It tells us how to
remove the leading deconvolution bias using only neighboring prefix coefficients.

For a function on the center lattice, define the centered second difference
\begin{equation}
 \Deltah^2H(\lambda)
 =\frac{H(\lambda+h_n)-2H(\lambda)+H(\lambda-h_n)}{h_n^2}. \label{eq:discrete-laplacian}
\end{equation}
Put
\begin{equation}
 c_n=\frac{3n+4}{72}h_n^2
\end{equation}
and define the corrected grid
\begin{equation}
 H_n^{[1]}(\lambda)=H_n(\lambda)+c_n\Deltah^2H_n(\lambda). \label{eq:corrected-grid}
\end{equation}
Every value in \eqref{eq:corrected-grid} is a rational linear combination of three normalized
iterated-prefix values.

\begin{theorem}[One-step deconvolution acceleration]
\label{thm:accelerated}
Uniformly on the center lattice,
\begin{equation}
 H_n^{[1]}(\lambda)-\Up(\lambda)
 =-\frac{n^2}{1152\,16^n}\Up^{(4)}(\lambda)
  +\smallo\!\left(\frac{n^2}{16^n}\right).                \label{eq:accelerated-profile}
\end{equation}
Consequently,
\begin{equation}
 \lim_{n\to\infty}\frac{16^n}{n^2}
 \sup_{\lambda\in h_n(\Z-g_n/2)}
 |H_n^{[1]}(\lambda)-\Up(\lambda)|
 =\frac89.                                                 \label{eq:accelerated-norm}
\end{equation}
\end{theorem}

\begin{proof}
For a smooth function,
\begin{equation}
 \Deltah^2f=f''+\frac{h_n^2}{12}f^{(4)}+\bigO(h_n^4).      \label{eq:second-difference-expansion}
\end{equation}
Write \eqref{eq:center-expansion} as
$H_n=\Up-c_n\Up''+d_n\Up^{(4)}+\bigO(n^3h_n^6)$, where
\[
 d_n=\left(\frac{15n+16}{43200}
          +\frac{(3n+4)^2}{10368}\right)h_n^4.
\]
Substitute this into \eqref{eq:corrected-grid} and use
\eqref{eq:second-difference-expansion}.  If the uniform remainder is denoted by $R_n$, then
$\|\Deltah^2R_n\|_\infty\le4h_n^{-2}\|R_n\|_\infty$; hence
$c_n\Deltah^2R_n=\bigO(n^4h_n^6)=\smallo(n^2h_n^4)$.  Thus the remainder remains negligible
at the claimed scale.  The second-derivative terms cancel, and the coefficient of
$\Up^{(4)}$ is
\begin{equation*}
 d_n-c_n^2+\frac{c_nh_n^2}{12}
 =-\frac{n^2}{1152}h_n^4+\bigO(nh_n^4).
\end{equation*}
This proves \eqref{eq:accelerated-profile}.  By \eqref{eq:sharp-derivatives},
$\|\Up^{(4)}\|_\infty=2^{10}=1024$, so the sharp norm constant is
$1024/1152=8/9$.
\end{proof}

Higher corrections follow by replacing the reciprocal-tail multiplier with successively
longer polynomials in the lattice Laplacian.  The all-orders expansion
\eqref{eq:all-orders} supplies the coefficients recursively.  To preserve the accelerated
rate in a continuous reconstruction, one should use interpolation of sufficiently high order:
piecewise-linear interpolation has an $\bigO(h_n^2)$ geometric error and would mask the
$\bigO(n^2h_n^4)$ corrected lattice error, whereas local cubic interpolation has
$\bigO(h_n^4)$ error and retains the latter as the dominant term.

\section{Conceptual synthesis}
\label{sec:synthesis}

\subsection{One row, three coordinate choices}

The normalized integer row
\begin{equation}
 m\longmapsto\frac{\sigma_{n+1}(m)}{2^{\binom n2}}        \label{eq:one-row}
\end{equation}
contains highly accurate samples, but the meaning of $m$ must be specified.

\begin{enumerate}
\item The intrinsic coordinate is the center
$\lambda_{n,m}=h_n(m-g_n/2)$.  At that coordinate the row differs from $\Up$ by a curvature
term of order $nh_n^2$.
\item Holding a center value constant across its width-$h_n$ cell changes the coordinate by
at most $h_n/2$.  The target has maximal slope $2$, so the histogram error is $h_n$.
\item Replacing the intrinsic coordinate by the naive dyadic coordinate implicit in
$m=\lfloor2^nx\rfloor$ overlooks the support deficit $r_n$.  It changes the coordinate by
roughly $nh_n/2$, so the Fabius error is roughly $nh_nF'(x)/2$.
\end{enumerate}
The three rates are therefore compatible, not contradictory.  They measure distinct
reconstructions of the same discrete data.

\subsection{Variance versus support radius}

The center expansion sees the omitted tail only through its even cumulants.  Its mean is zero,
so the first nonzero effect is its variance $v_n\asymp nh_n^2$.  The floor-coordinate drift,
by contrast, sees the extreme support loss $r_n\asymp nh_n$.  This is why a centered analytic
argument is exponentially more accurate than a support-aligned but miscentered index rule.

This distinction is common in lattice approximation.  Weak convergence controls averages,
and moment matching controls smooth centered observables, but an endpoint-derived coordinate
can retain a worst-case support defect long after the bulk distribution has become extremely
close.

\subsection{Why the \texorpdfstring{$n$}{n} is structural}

The degree
\begin{equation}
 g_n=2^{n+1}-n-2
\end{equation}
is shorter than the naive value $2^{n+1}$ by exactly $n+2$.  After scaling by
$2^{n+1}$ and centering, half of this deficit appears at each endpoint.  Hence the factor $n$
in the floor rate is already encoded in the elementary degree formula for $P_n$.  No delicate
property of $F$ is needed to predict its presence; the analytic work is needed to turn the
horizontal deficit into the sharp vertical profile $F'$ and to control the smaller terms.

\subsection{A practical recommendation}

For numerical reconstruction from a computed prefix row:
\begin{itemize}
\item do not use $m/2^n$ as the physical coordinate;
\item attach each value to $\lambda_{n,m}$ from \eqref{eq:center};
\item use at least linear interpolation rather than a histogram;
\item if another factor of roughly $4^n/n$ is valuable, apply the three-point correction
      \eqref{eq:corrected-grid} and use cubic interpolation.
\end{itemize}
The coefficients remain exact rationals until the final interpolation or evaluation stage.

\section{A Lean formalization roadmap}
\label{sec:lean-roadmap}

The new proof separates naturally into reusable lemmas.  A formalization need not begin with
the full asymptotic theorem.

\begin{enumerate}
\item \textbf{Shifted-lattice Fourier inversion.}  Formalize \cref{lem:lattice-inversion} for
a finitely supported measure on $a+h\Z$.  This is a finite-sum orthogonality calculation.
\item \textbf{Rademacher realization.}  Refine the existing finite-measure bridge to the
triangular Rademacher array \eqref{eq:Xn}; obtain \eqref{eq:finite-cosine} and the independent
tail decomposition \eqref{eq:Tn}.
\item \textbf{Exact geometric tails.}  Add the closed forms \eqref{eq:rn}, \eqref{eq:vn},
\eqref{eq:alpha-n}, and \eqref{eq:beta-n}.  These reduce to finite manipulations of geometric
series.
\item \textbf{Uniform $\log\sec$ remainder.}  Prove a sixth-order Taylor bound on
$[-1/4,1/4]$ and use it scale by scale.  The low-frequency cutoff can be chosen as
$2^{n/4}$ or replaced by any convenient dyadic integer bound.
\item \textbf{Rapid Fourier decay.}  Package repeated integration by parts for compactly
supported $C^\infty$ functions and instantiate it with the already formalized sharp derivative
bounds.
\item \textbf{Center expansion.}  Combine the preceding lemmas with Fourier derivative
identities.  The first-order version \eqref{eq:center-leading} is substantially shorter than
the fourth-order theorem and already proves all three basic rates.
\item \textbf{Exact drift.}  Formalize \eqref{eq:lambda-drift} by integer-floor arithmetic.
Then \eqref{eq:floor-first} is a direct Taylor estimate.
\item \textbf{Midpoint identity.}  Prove \cref{thm:midpoint-exact} independently of Fourier
analysis from palindromicity and the central coefficient plateau.  It is an especially useful
regression test for index conventions.
\end{enumerate}

The formal theorem statements should distinguish three domains: values at the center lattice,
the step function on all reals, and the floor sample on $[0,1]$.  Combining them into one
statement would obscure exactly the coordinate issue responsible for the different rates.

\section{Conclusion}
\label{sec:conclusion}

Iterated summation of the signed Thue--Morse sequence does more than produce a sequence that
happens to resemble the Fabius function.  After the correct one-order shift, its finite row is
an exact lattice deconvolution of $\Up$.  The omitted part is an explicit triangular Bernoulli
tail with support radius $(n+2)2^{-n-1}$ and variance $(3n+4)/(36\,4^n)$.  Those two exact
numbers govern two different approximation regimes.

At the natural centers, the mean-zero tail leaves a variance-scale curvature bias, yielding
the sharp rate $n/(3\,4^n)$.  Holding the values constant across cells introduces the larger
sharp rate $2^{-n}$.  Reading the same row at the naive dyadic floor leaves the full support
deficit as a coordinate drift and yields the still larger sharp rate $n2^{-n}$.  The latter
has the uniform error profile $F'$ and an exact midpoint formula.

The main conceptual correction is therefore simple but consequential: the integer index is
not itself the continuum coordinate.  Once the genuine center map is retained, the
Thue--Morse prefix construction is far more accurate than its original pointwise formulation
suggests, and its error admits a complete, explicitly computable differential expansion.

\appendix

\section{Elementary geometric-tail formulas}
\label{app:tails}

For $|q|<1$ and $n\ge0$,
\begin{equation}
 \sum_{\ell=n+1}^{\infty}\ell q^\ell
 =q\frac{\dd}{\dd q}\left(\frac{q^{n+1}}{1-q}\right)
 =q^{n+1}\frac{(n+1)-nq}{(1-q)^2}.                        \label{eq:geometric-tail}
\end{equation}
The specializations used in the text are
\begin{align}
 \sum_{\ell>n}\ell2^{-\ell}&=(n+2)2^{-n},\\
 \sum_{\ell>n}\ell4^{-\ell}&=\frac{3n+4}{9\,4^n},\\
 \sum_{\ell>n}\ell16^{-\ell}&=\frac{15n+16}{225\,16^n},\\
 \sum_{\ell>n}\ell64^{-\ell}&=\frac{63n+64}{3969\,64^n}.
\end{align}
These identities make every coefficient in the all-orders expansion elementary: no
unevaluated infinite tail sum is necessary.

\section{A short exact-arithmetic algorithm}
\label{app:algorithm}

The coefficient row $[z^m]P_n$ can be built without polynomial multiplication.  Multiplication
by $1+\cdots+z^{2^i-1}$ is a sliding-window sum.  The following pseudocode maintains exact
integers and costs $\bigO(2^n)$ arithmetic operations per level.

\begin{lstlisting}[style=pseudo]
coeff := [1]
for i := 1,...,n:
    width := 2^i
    nextLength := length(coeff) + width - 1
    next := array(nextLength, 0)
    window := 0
    for m := 0,...,nextLength-1:
        if m < length(coeff):
            window := window + coeff[m]
        if m >= width:
            window := window - coeff[m-width]
        next[m] := window
    coeff := next

# Normalized prefix sample:
H_n(m) := coeff[m] / 2^(n(n-1)/2)
# Physical coordinate:
lambda_n(m) := (2m - (2^(n+1)-n-2)) / 2^(n+1)
\end{lstlisting}

The midpoint identity \eqref{eq:midpoint-exact}, symmetry
$\mathtt{coeff[m]}=\mathtt{coeff[g_n-m]}$, the total mass
$2^{\binom{n+1}{2}}$, and the central maximum $2^{\binom n2}$ provide independent exact
checks of an implementation.

\begin{thebibliography}{99}

\bibitem{fabius1966}
J.~Fabius,
\newblock A probabilistic example of a nowhere analytic $C^\infty$-function,
\newblock \emph{Zeitschrift f\"ur Wahrscheinlichkeitstheorie und Verwandte Gebiete}
  \textbf{5} (1966), 173--174.
\newblock DOI: \nolinkurl{10.1007/BF00536652}.

\bibitem{arias1982}
J.~Arias de Reyna,
\newblock Definici\'on y estudio de una funci\'on indefinidamente diferenciable de soporte
  compacto,
\newblock \emph{Revista de la Real Academia de Ciencias Exactas, F\'isicas y Naturales de
  Madrid} \textbf{76} (1982), no.~1, 21--38.
\newblock English translation: \emph{An infinitely differentiable function with compact
support: Definition and properties}, arXiv:\nolinkurl{1702.05442} (2017).

\bibitem{allouche2003}
J.-P.~Allouche and J.~Shallit,
\newblock \emph{Automatic Sequences: Theory, Applications, Generalizations},
\newblock Cambridge University Press, 2003.

\bibitem{reshetnikovExposition}
V.~Reshetnikov,
\newblock \emph{The Fabius Function and Rvachev's Up-Function: Exact arithmetic,
probability, Fourier analysis, and corrected sharp asymptotics},
\newblock human-readable synthesis in the \texttt{ProveIt} repository, snapshot
25 August 2026.
\newblock \url{https://github.com/VladimirReshetnikov/ProveIt/blob/e4aa81a0b6b7e4f9943c3a5789c4a781ee120be4/Analysis/FabiusFunction/docs/Fabius_Function_and_Rvachev_Up/Fabius_Function_and_Rvachev_Up.tex}.

\bibitem{proveit}
V.~Reshetnikov,
\newblock \emph{ProveIt: Analysis/FabiusFunction},
\newblock Lean~4 formal development, repository snapshot
\path{e4aa81a0b6b7e4f9943c3a5789c4a781ee120be4}, 25 August 2026.
\newblock \url{https://github.com/VladimirReshetnikov/ProveIt/tree/e4aa81a0b6b7e4f9943c3a5789c4a781ee120be4/Analysis/FabiusFunction}.

\end{thebibliography}

\end{document}
